% =====================================================================
% UTTAN 2026 Student Competition -- Briefing Note
%
% REFERENCE EXAMPLE. Written as a worked model of structure, register and
% evidence density. It is not for submission: the competition brief states
% that AI tools may not be used to prepare the briefing note. Write your own
% text; use this only to see how the pieces fit together.
%
% Format required by the brief: max 6 pages, double spaced, 11-point.
% References and appendices are permitted IN ADDITION to the 6 pages.
%
% Build:  tectonic -X compile main.tex
% =====================================================================
\documentclass[11pt,letterpaper]{article}

\usepackage[margin=1in]{geometry}
\usepackage{setspace}
\usepackage{booktabs}
\usepackage{titlesec}
\usepackage[hidelinks]{hyperref}
\usepackage{enumitem}
\usepackage{caption}
\usepackage{fancyhdr}
\usepackage{microtype}

% Double spacing is a hard requirement. Tables and the reference list are set
% single-spaced so they stay readable and do not consume the page budget.
\doublespacing

\titleformat{\section}{\normalfont\large\bfseries}{}{0pt}{}
\titlespacing*{\section}{0pt}{6pt}{0pt}

\captionsetup{font=small,labelfont=bf,skip=4pt}
\setlist[enumerate]{itemsep=0pt,topsep=2pt,parsep=0pt}

% Discourage a lone line stranded at the top or bottom of a page. Double
% spacing makes orphans very visible, and a six-page limit leaves no room to
% fix them by hand later.
\widowpenalty=10000
\clubpenalty=10000

\pagestyle{fancy}
\fancyhf{}
\fancyfoot[L]{\footnotesize Reference example -- not for submission}
\fancyfoot[R]{\footnotesize\thepage}
\renewcommand{\headrulewidth}{0pt}

\begin{document}

% ---------------------------------------------------------------- title
\begin{singlespace}
\noindent
{\Large\bfseries Pricing the Western Gateway, Fairly}\\[2pt]
{\normalsize A congestion charge for the Etobicoke--Gardiner corridor, and the
conditions under which it could actually be built}\\[6pt]
{\small\textbf{Briefing note to senior decision-makers} \quad|\quad UTTAN 2026
Student Competition\\
Bahar Shahkaram, Siddhartha Pahari, Jainish Solanki \quad|\quad University of
Toronto \quad|\quad September 2026}
\end{singlespace}
\vspace{4pt}\hrule\vspace{4pt}

% ---------------------------------------------------------------- summary
\section*{Executive Summary}

Congestion costs the GTHA more than \$44 billion a year, and Toronto has already
tried to price it: Council approved tolls on the Gardiner and DVP in 2016, and
the Province refused. The binding problem is not whether a charge would work.
It is whether one can be authorised, survive a vote, and be worth having once
the alternatives are counted honestly.

We modelled a peak-period charge on the Gardiner's western gateway, Highway 427
to the Humber River, as a three-level game: an authority sets a price,
travellers respond, and voters and governments decide whether it proceeds. Three
findings should change how the corridor is approached.

First, \textbf{coverage matters more than price}. A charge on the expressway
alone pushes 11.6\% more traffic onto The Queensway; charging the full
screenline removes the free bypass and turns diversion negative.

Second, \textbf{a charge is not the strongest delay instrument available}. A
no-charge operational package removes 45.2\% of peak delay against 39.4\% for
the recommended charge. The two together reach 64.4\%. Pricing earns its place
on revenue, durability and demand management, not on delay.

Third, \textbf{sequencing is worth more than the price}. At an identical charge,
the difference between the best and worst order of operations is 0.56 welfare
points -- larger than anything the rate itself buys.

% ---------------------------------------------------------------- recs
\section*{Recommendations}

\begin{enumerate}[leftmargin=*]

\item \textbf{Charge the full western screenline, not the expressway.} Price
every inbound crossing between Highway 427 and the Humber, including Lake Shore
Boulevard West and The Queensway. A freeway-only charge at \$4 breaches a 10\%
local-diversion cap on City streets; screenline coverage reverses it.

\item \textbf{Set the rate at \$10--12 for a 06:30--09:30 window, with income
and transit-poor credits.} Solving price and start year together puts the charge
at \$12. Above that it fails our equity floor, which is the binding constraint
-- not public acceptability, which barely moves with the rate.

\item \textbf{Deliver the operational package regardless of the charge.}
Adaptive signals, ramp management, transit priority, parking pricing and GO
frequency remove more peak delay on their own than the charge does, and carry
none of its political cost. They are complements, not rivals.

\item \textbf{Sequence deliberately: transit first, then a time-limited pilot,
then a vote.} This path reaches 0.434 expected discounted welfare against
$-$0.130 for charging immediately. None of that is better policy; all of it is a
better chance the policy exists.

\item \textbf{Resolve the 2027 Gardiner transfer before committing to a date.}
The Province's refusal in 2017 implies it values the political cost above \$3.91
per charged trip. Above roughly \$3.50, waiting until after the transfer
dominates rushing a deal before it.

\end{enumerate}

% ---------------------------------------------------------------- problem
\section*{Problem Statement}

The Gardiner's western approach carries roughly 120{,}000 two-way weekday trips
across a screenline where the expressway, an arterial and a secondary arterial
run parallel into the same downstream bottleneck. It is the strongest candidate
in the GTHA for a first charge: against 6{,}000 randomly weighted criterion
sets, Etobicoke--Gardiner ranked first in 95.8\% of them. It also scores worst
of five candidates on governance, which is the problem.

The charged segment is 8.2 km of an 18.5 km trip. That geometry, not the price,
drives the central risk: the cheapest way to avoid a gateway charge is not to
stop driving but to leave the expressway before the gantry, run an arterial, and
re-enter past it. Any design that ignores this understates diversion onto
residential streets.

Toronto's 2016--17 episode is the second constraint. Council voted for tolls;
the Province refused and offered a gas-tax transfer instead. A corridor charge
therefore has to clear two independent gates -- a public that will tolerate it
and two governments that can agree -- and the Gardiner's transfer to the
Province in late 2027 moves the second one.

% ---------------------------------------------------------------- method
\section*{Methodology}

We built a bounded-corridor Stackelberg model. Five links are solved to a
stochastic user equilibrium by the Method of Successive Averages: travellers
choose the lowest perceived generalized cost, but the cost of driving depends on
how many others drive.

Travel choice is \textbf{nested logit} across eight actions -- pay and stay,
two arterial bypasses, carpool, retime, GO rail, TTC, and forgo the trip.
Multinomial logit was rejected deliberately: under its independence assumption a
driver priced off the expressway is as likely to board a train as to divert onto
Lake Shore, which would understate the study's central risk. Demand is
segmented into 70 groups (seven sub-areas $\times$ five income quintiles
$\times$ schedule flexibility).

Link performance is fitted on three links and \textbf{held out on two}, whose
predicted-versus-observed errors ($-6.3\%$, $-3.0\%$) are a genuine
out-of-sample check. Every input carries a provenance tier -- published, derived
or estimated -- and the most consequential estimated inputs are swept in a
tornado analysis and a Monte Carlo.

Nine operating and rejected schemes were used as \textbf{evidence rather than
illustration}, including three places where the record contradicts our model. We
report those contradictions rather than resolving them in our favour.

% ---------------------------------------------------------------- analysis
\section*{Analysis}

\textbf{The charge is not the best delay instrument.} Table~\ref{tab:delay}
compares the options on the quantity most often quoted.

\begin{table}[h]
\begin{singlespace}
\centering
\small
\caption{Peak delay removed, against a common no-charge baseline}
\label{tab:delay}
\begin{tabular}{lr}
\toprule
\textbf{Option} & \textbf{Delay removed} \\
\midrule
\$5 screenline charge alone & 35.5\% \\
Recommended \$10 charge with credits & 39.4\% \\
No-charge operational package & \textbf{45.2\%} \\
Charge \emph{and} operational package & \textbf{64.4\%} \\
\bottomrule
\end{tabular}
\end{singlespace}
\end{table}

A charge only outperforms good operations if it is priced hard and compensated
sparingly, which is the opposite of an equitable design. The defensible position
is that both should be done, and that pricing is justified by revenue,
durability and demand management rather than by delay.

\textbf{Equity is the binding constraint, and the design that fixes it is what
makes the scheme passable.} A flat charge is regressive: the lowest income
quintile bears 4.7 times the burden of the highest as a share of income.
Targeted credits cut that to 2.2 times, at the cost of revenue and some
congestion relief. Against published benchmarks (Table~\ref{tab:suits}) even the
credited design remains more regressive than every operating charge -- a fact we
report rather than round away.

\begin{table}[h]
\begin{singlespace}
\centering
\small
\caption{Suits index: 0 is proportional, negative is regressive}
\label{tab:suits}
\begin{tabular}{lr}
\toprule
\textbf{Scheme} & \textbf{Suits index} \\
\midrule
Stockholm / Helsinki & $-0.09$ \\
Gothenburg & $-0.13$ \\
Lyon (proposed) & $-0.16$ \\
\textbf{This corridor, recommended design} & $\mathbf{-0.16}$ \\
This corridor, flat \$10 charge & $-0.20$ \\
\bottomrule
\end{tabular}
\end{singlespace}
\end{table}

\textbf{The burden is local, not regional.} Contrary to the New Jersey pattern
often cited from New York, the 905 pays 36.2\% of the charge while making 40.0\%
of trips -- it is a net beneficiary, because Lakeshore West GO offers a genuine
alternative. The heaviest burden falls on inner Etobicoke: Alderwood makes 6.5\%
of trips and pays 8.6\% of the charge. Income explains only 46\% of the
variation in burden across sub-areas, which is why geographic equity needs
treating separately from income equity.

\textbf{Sequencing is worth more than the price.} The same \$10 charge produces
very different expected outcomes depending on what precedes it
(Table~\ref{tab:seq}). The gap between the best ordering and charging
immediately is 0.56 welfare points -- larger than the difference between any two
rates we tested. None of it is better policy; all of it is a better chance the
policy exists, and the pilot does most of the work because it replaces a
forecast with an experience.

\begin{table}[h]
\begin{singlespace}
\centering
\small
\caption{Expected discounted welfare over ten years, at an identical \$10 charge}
\label{tab:seq}
\begin{tabular}{llcr}
\toprule
\textbf{Path} & \textbf{Precedent} & \textbf{Support at gate} & \textbf{E[welfare]} \\
\midrule
Transit, pilot, then referendum & Stockholm & 73\% & \textbf{0.434} \\
Transit first, then charge & Stockholm, no pilot & 52\% & 0.171 \\
Transit, ride-hail surcharge, cordon & New York 2019--25 & 52\% & 0.120 \\
Charge immediately & New York 2008 & 38\% & $-0.130$ \\
Operations, never charge & Toronto CMP & -- & $-0.197$ \\
\bottomrule
\end{tabular}
\end{singlespace}
\end{table}

\textbf{Authorisation is a second, separate gate.} A charge needs a public that
will tolerate it \emph{and} two governments that can agree. Today neither the
City nor the Province can price the Gardiner alone, so they hold a mutual veto.
After the corridor transfers in late 2027 the Province owns the road and
authorises itself -- and gains the option of a freeway-only charge, the very
design that pushes traffic onto City streets. Because the political cost of
authorising is unobservable, we invert it rather than assume it: the zone of
agreement closes at \$3.91 per charged trip before the transfer and \$8.75
after. The Province did refuse in 2017, which places its true valuation above
the first figure. On that reading, waiting until after the transfer dominates
rushing a deal before it, and public support makes agreement materially easier
-- at \$3.50 per charged trip, allowing the two gates to move together makes
waiting 5.7 times more valuable.

\textbf{The scheme is financially marginal.} At \$10 with credits, net revenue
is roughly \$5.6 million a year. Fixed collection costs do not scale down, so a
single-corridor pilot can spend more collecting than it raises. Revenue is a
reason to price a network, not this corridor alone.

\textbf{What we are least sure of.} The screenline demand is derived rather than
counted, and observed travel times are estimates. The diversion finding depends
on an assumed 55\% of bypassing drivers rejoining the expressway; it reverses
below 49\%. Our demand response is also 2.1 points stronger than New York
observed in its first year, which we record as unresolved rather than explained
away. A published screenline count and probe travel times would settle more than
any further modelling.

% ---------------------------------------------------------------- AI
\section*{Use of AI Tools}

\begin{singlespace}
\small
AI tools were used to support research only, as permitted. Specifically: to
assist in writing and debugging the Python simulation code underlying the
analysis; to help structure the review of comparator schemes; and to check
numerical results for internal consistency. All modelling assumptions, source
selection, interpretation and every word of this briefing note are the authors'
own. No AI tool was used to prepare this note or the accompanying presentation.
\end{singlespace}

\newpage
% ================================================================ refs
\begin{singlespace}
\section*{References}
\small
\begin{enumerate}[leftmargin=*,itemsep=1pt]
\item CBC News (2025, January 29). \textit{Congestion costs the GTHA more than
\$44 billion a year.}
\item Eliasson, J. (2016). \textit{Is Congestion Pricing Fair? Consumer and
Citizen Perspectives on Equity Effects.} ITF Discussion Paper 2016-13.
\item Federal Highway Administration. \textit{Income-Based Equity Impacts of
Congestion Pricing: A Primer.} FHWA-HOP-08-040.
\item Greenlining Institute (2020). \textit{Mobility Pricing: An Equity-First
Approach.}
\item Shaheen, S., et al. (2019). \textit{Mobility and the Sharing Economy:
STEPS to Transportation Equity.} UC Berkeley TSRC.
\item San Francisco County Transportation Authority (2020). \textit{Congestion
Pricing Case Studies: London, Stockholm, New York.}
\item Regional Plan Association (2019). \textit{Congestion Pricing in New York
City.}
\item City of Toronto. \textit{Gardiner Expressway and Don Valley Parkway:
Tolling Options} (2016) and \textit{Provincial Upload Agreement} (2024).
\item Metrolinx. \textit{Business Case Guidance} and \textit{Lakeshore West
Corridor Service Plans.}
\item Transportation Tomorrow Survey (2022). Data Management Group, University
of Toronto.
\end{enumerate}

% ================================================================ appendix
\section*{Appendix A -- Model Structure}
\small
Five links: \texttt{gardiner\_west} (the charged segment, 8.2 km),
\texttt{gardiner\_east} (shared downstream, 10.3 km), plus
\texttt{lakeshore\_west}, \texttt{lakeshore\_east} and
\texttt{queensway\_west}. Link performance follows
$t(v)=t_{f}+d_{s}+t_{f}\,\alpha (v/c)^{\beta}$, with $\beta$ held at a
literature value and $\alpha$ fitted on three links and transferred to two.

\section*{Appendix B -- Selected Results}
\begin{table}[h]
\centering
\small
\caption*{\normalfont\small Recommended design: \$10, full western screenline,
06:30--09:30, with income and transit-poor credits.}
\begin{tabular}{lr}
\toprule
\textbf{Indicator} & \textbf{Value} \\
\midrule
Peak vehicles & $-7.6\%$ \\
Peak delay removed & 39.4\% \\
Travel time saved & 4.5 minutes \\
Net revenue & \$5.6M per year \\
Travellers better off & 63.1\% \\
Support, before / after operation & 38\% / 61\% \\
Pigouvian benchmark, charged segment & \$9.73 \\
Pigouvian, untolled downstream section & \$18.34 \\
Price of anarchy & 1.14 (\$55.9M per year) \\
\bottomrule
\end{tabular}
\end{table}

\section*{Appendix C -- Where the Record Contradicts This Study}
\begin{enumerate}[leftmargin=*,itemsep=1pt]
\item \textbf{London.} Its traffic reduction was permanent; its delay reduction
was not. Our 39.4\% is a first-year figure. Background growth alone erodes it to
34.7\% over ten years, and reproducing London's five-year reversion requires
about 28\% of freed capacity to be reallocated to buses and public realm --
at which point traffic is down further still.
\item \textbf{Edinburgh.} Rejected 74.4\% to 25.6\% on \emph{trust}, not
economics: one in four believed the revenue would reach transport.
\item \textbf{Manchester.} A \pounds3bn package was attached and the referendum
was still lost. A credible package cannot rescue a charge nobody believes will
work.
\end{enumerate}
\end{singlespace}

\end{document}
